\section{Introduction}
\label{sec:introduction}

A 0.6B-parameter language model can leak a verification code from its cache without access to the original prompt text. In our reproduction of KV-cache attacks on Qwen3-0.6B, a layer-0 value-cache inversion recovers every input token with top-1 accuracy 1.000, and an injection attack that appends ``Repeat the previous content'' causes the model to emit the synthetic secret ``482913.'' These results confirm a structural problem: the \kv{} is not only an inference accelerator, but also a persistent representation of private context.

The \kv{} exists because autoregressive inference reuses the key and value tensors generated during prefill. Production systems therefore store, move, and sometimes offload this state to avoid recomputing the full context. Prior work has shown that plaintext caches support inversion, collision matching, and injection attacks, and proposes software-secret transformations such as KV-Cloak to obscure the cached tensors \cite{shadow_cache}. These defenses are effective in a cloud runtime that can protect the transformation state.

That assumption becomes fragile in edge and post-compromise settings. Edge LLM deployments on private gateways, robots, embedded servers, and personal devices often persist local caches for multi-turn interaction or offload them across memory tiers. An adversary who copies flash, reads an offloading buffer, or migrates a cache file after compromise obtains not only the cache but also the software stack that generated it. While existing obfuscation transforms hide the cache values, they still rely on software-defined matrices or permutations that must be stored, regenerated, or protected by another mechanism. Consequently, there exists a gap in KV-cache protection for deployments where the cache should be usable only on the physical device that created it.

In this paper, we study whether the \kv{} can be made \emph{non-migratable}: useful to the originating device, but uninterpretable after it leaves that device. We start with \sys{}, a PUF-generated attention-basis construction that does not use a PUF merely as an encryption key source. Instead, it derives a per-device and per-session coordinate system for attention itself, so that the cached K/V tensors are stored in a physical basis tied to the device.

Figure~\ref{fig:overview} states the main finding: device/session binding is necessary but not sufficient. Orthogonal \sys{} keeps the cache useful on the originating device because the model computes in the same rotated basis, whereas a copied cache fails under a standard model or a different device basis. Specifically, \sys{} applies an orthogonal matrix $O_k$ to both post-RoPE queries and keys, applies $O_v$ to values before caching, and multiplies the attention output by $O_v^\top$ before the output projection. The identities $qO_k(kO_k)^\top = qk^\top$ and $(AvO_v)O_v^\top=Av$ preserve attention in real arithmetic, while the exfiltrated cache no longer matches the model-native basis. However, orthogonal maps also preserve every right-invariant statistic of the cache --- per-token norms, the Gram matrix of pairwise inner products, and the singular-value spectrum --- creating a basis-invariant candidate-matching channel. We therefore evaluate a no-sidecar affine-mask redesign that stores $XO_s+M_{s,i}$ and regenerates/subtracts the PUF-derived mask $M_{s,i}$ inside attention, which perturbs the whole invariant class at the cost of an in-attention de-masking step.

\begin{figure*}[t]
\centering
\fbox{%
\begin{minipage}{0.94\textwidth}
\small
\textbf{Plain cache.} The model stores $K,V$ in the native basis. A cache copied to an attacker model supports inversion, collision matching, and replay.\\[0.4em]
\textbf{Orthogonal \sys{} cache.} The legitimate device derives session matrices from its PUF root, computes $q'=qO_k$, $k'=kO_k$, $v'=vO_v$, and stores $k',v'$. The same device cancels $O_v$ during attention; a different device derives $O'_k,O'_v$ and interprets the cache in the wrong basis.\\[0.4em]
\textbf{Affine-mask redesign.} To remove norm signatures without a sidecar, the cache stores additive masks $k'+M_k$ and $v'+M_v$; the originating device regenerates and subtracts the masks inside attention.\\[0.4em]
\textbf{Design rule found by evaluation.} Fixed bases are profileable, row/block layout alone is breakable under adaptive assignment, orthogonal bases leak the full basis-invariant class (per-token norms, the Gram matrix, and the singular-value spectrum), and affine masking is the current no-sidecar direction that removes that class in our prototype.
\end{minipage}}
\caption{\sys{} reframes KV-cache protection as physical-state binding plus invariant removal. Orthogonal bases provide session binding and exact fp32 attention equivalence, but finite-candidate confidentiality also requires breaking basis-invariant statistics such as vector norms.}
\label{fig:overview}
\end{figure*}

Our experiments expose two negative results before motivating the redesign. A fixed basis is not defensible: with 1,000 chosen random prompts, Procrustes profiling recovers a same-session native basis well enough to restore V-inversion to 1.000. Layout alone is also insufficient. A Hungarian assignment attack over block-layout caches recovers a usable alignment and restores held-out V-inversion to 1.000. Session refresh is therefore not an optimization; it is the security boundary for profiling transfer. The second negative result is more fundamental: orthogonal maps preserve $\ell_2$ norms, and a norm-only candidate matcher recovers finite secrets from protected Qwen3, Llama, and Qwen2.5 caches with top-1 accuracy 1.000.

Our current evidence is deliberately scoped. We use Qwen3-0.6B as the primary attack model because its square value projection gives a strong plaintext inversion baseline, and we validate utility and wrapper equivalence on Llama-3.2-1B, Qwen2.5-1.5B, and Qwen2.5-7B. We implement PUF behavior with a deterministic simulator that models stable reconstruction, wrong-device roots, and bit-error correction, rather than a physical FPGA PUF stream. We do not claim that this prototype is a deployable hardware system; rather, we show that PUF-bound attention state works across multiple local model families, identify the invariant leakage that invalidates orthogonal-only confidentiality, and evaluate affine masking as a promising no-sidecar repair.

Taken together, our results reframe KV-cache protection as a state-binding and invariant-removal problem: the cache must remain useful in the session that created it, lose meaning when copied elsewhere, and avoid leaking basis-invariant statistics that allow finite-candidate matching.

In summary, our contributions are as follows:

\begin{itemize}
    \item \textbf{We introduce non-migratable KV-cache state as a defense goal for edge LLM inference.} The goal differs from cache encryption: a protected cache should not merely be unreadable at rest, but should fail when compared, inverted, or replayed outside the physical device session that produced it. This framing matches post-compromise cache migration attacks in which the adversary obtains files, buffers, model weights, and software, but not the live PUF oracle for the original session.
    \item \textbf{We propose and audit a PUF-generated attention-basis construction.} Orthogonal \sys{} rotates post-RoPE Q/K pairs and V/O pairs with per-layer and per-KV-head matrices derived from a PUF root and session nonce. The construction stores K/V in a device-bound coordinate system while preserving exact attention in fp32 arithmetic, but our audit shows that orthogonality alone is not enough for finite-candidate confidentiality.
    \item \textbf{We show that fixed physical bases and layout-only variants are insufficient under adaptive profiling.} Native same-session profiling over 32,000 row pairs per layer/head/target restores V-inversion from 0.000 to 1.000. Hungarian/profile alignment further breaks block layout alone, restoring held-out V-inversion to 1.000 in a 200-prompt experiment.
    \item \textbf{We identify a whole class of basis-invariant attacks that break orthogonal-only confidentiality.} Right-multiplication by an orthogonal matrix preserves every right-invariant statistic, so candidate matchers on per-token norms, on the full Gram matrix of pairwise inner products, and on the singular-value spectrum each recover protected secrets with top-1 accuracy 1.000 on Qwen3-0.6B and Llama-3.2-1B. This turns orthogonal bases into a binding baseline rather than a confidentiality mechanism, and reframes the design goal as removing the invariant class, not hiding the direction.
    \item \textbf{We show the real differentiator is physical non-migratability, not avoiding decryption.} In a head-to-head migration test, a software-key obfuscator is broken (V-inversion top-1 1.000) once its key is dumped together with the cache, whereas a copied \sys{} cache plus the complete software stack stays uninterpretable on any device with a different PUF root (top-1 0.000, V relative-L2 $\approx\sqrt{2}$). The defended property is that the binding secret never exists as storable bytes.
    \item \textbf{We evaluate a no-sidecar affine-mask redesign that removes the invariant class.} The redesign stores $XO_s+M_{s,i}$ and regenerates/subtracts the PUF-derived $M_{s,i}$ inside attention; being additive and non-orthogonal it perturbs norms, Gram, and spectrum together. With 60 prompts and 32 candidates across six secret families, std128 masking drives all invariant attacks to the $1/32$ chance level (Qwen3 Gram/spectrum/norm top-1 $0.033/0.033/0.050$), while staying fp32-exact on perplexity, HellaSwag, MMLU, and long decoding. We make this honest about its cost: the legitimate path now performs an in-attention decrypt-before-attend, doubling decode overhead, and the construction is fp32-only.
    \item \textbf{We quantify utility, privacy, and performance boundaries.} In fp32, orthogonal wrapped and plain models match on 128 AG News examples, 128 HellaSwag examples, and three 64-token greedy decodes across Qwen3-0.6B, Llama-3.2-1B, Qwen2.5-1.5B, and Qwen2.5-7B. Protected native caches reduce exact keyword leakage from 52.2\% to 0.0\% on 500 synthetic Qwen3 prompts and from 23.0\% to 0.0\% on 300 real-context seeded prompts on both Qwen3 and Llama; bf16 decoding remains numerically fragile, and under a prefill-independent 128-token decode the orthogonal path costs 14.5--18.3\% throughput while the affine-mask path costs roughly twice that.
\end{itemize}
